\documentclass{article}
\usepackage{ctex}
\usepackage{graphicx}
\usepackage{multicol}
\usepackage{caption}
\usepackage{blindtext}
\usepackage{cuted} % cuted宏包提供strip环境，用于解决figure*环境中的图片在双栏模式下无法在首页被插入的问题
%! 注意，在strip环境中需要使用caption宏包提供的"\captionof{float type}{heading}"命令来设置标题
% \usepackage{dblfloatfix} % dblfloatfix宏包用于解决figure*环境中的图片在双栏模式下无法在底部被插入的问题，但是仍然无法解决图片无法在首页被插入的问题
%! 注意，dblfloatfix宏包和stfloats宏包不兼容，只能加载其中一个
%* 但是dblfloatfix宏包可以和nidanfloat宏包兼容
\usepackage{stfloats} % stfloats宏包可以同时解决上述两个问题
% \usepackage{nidanfloat} % nidanfloat宏包也可以同时解决上述两个问题
%! 注意，加载nidanfloat宏包会导致通过cuted宏包提供的strip环境中的图片消失，但编译不会报错

\setlength{\columnseprule}{.4pt}

\begin{document}
% （一）在双栏模式下使用figure环境插入图片
\twocolumn
    \textbf{第一部分：figure环境}

    这是一段测试文字……这是一段测试文字……这是一段测试文字……这是一段测试文字……这是一段测试文字……这是一段测试文字……这是一段测试文字……这是一段测试文字……这是一段测试文字……这是一段测试文字……这是一段测试文字……这是一段测试文字……这是一段测试文字……

    \newpage 
    
    \begin{figure}[htbp]
        \includegraphics[width = .1\textwidth]{doraemon1.jpg}
        \caption{这是用figure环境在双栏模式下插入的图片（一）}
    \end{figure}
    这是一段测试文字……这是一段测试文字……这是一段测试文字……

    \begin{figure}[htbp]
        \includegraphics[width = .1\textwidth]{doraemon1.jpg}
        \caption{这是用figure环境在双栏模式下插入的图片（二）}
    \end{figure}
    这是一段测试文字……这是一段测试文字……
    \begin{figure}[htbp]
        \includegraphics[width = .1\textwidth]{doraemon1.jpg}
        \caption{这是用figure环境在双栏模式下插入的图片（三）}
    \end{figure}
    这是一段测试文字……这是一段测试文字……这是一段测试文字……

    %* 可以发现，在双栏模式下使用figure环境插入图片和在单栏模式下没有什么区别，只不过插入的区域从单栏变成了双栏，插入的图片可以在双栏内部浮动

% （二）在双栏模式下使用figure*环境插入图片
%! 由于加载了stfloats宏包，因此打印出的实际上是stfloats宏包的效果，如果想要查看在没有加载宏包的情况下，在双栏模式下使用figure*环境插入图片的效果，可以取消加载stfloats宏包
    \twocolumn
    %* 继续使用一次"\twocolumn[text]"命令，效果是另起一页继续双栏模式
    \textbf{第二部分：figure*环境}

    \begin{figure*}[htbp]
        \includegraphics[width = .1\textwidth]{doraemon1.jpg}
        \caption{这是用figure*环境在双栏模式下插入的图片（一）}
    \end{figure*}
    这是一段测试文字……这是一段测试文字……这是一段测试文字……

    \begin{figure*}[htbp]
        \includegraphics[width = .1\textwidth]{doraemon1.jpg}
        \caption{这是用figure*环境在双栏模式下插入的图片（二）}
    \end{figure*}
    这是一段测试文字……这是一段测试文字……
    \begin{figure*}[htbp]
        \includegraphics[width = .1\textwidth]{doraemon1.jpg}
        \caption{这是用figure*环境在双栏模式下插入的图片（三）}
    \end{figure*}
    这是一段测试文字……这是一段测试文字……这是一段测试文字……

    \blindtext[7]

    %* 以上做了两步操作：（一）将上一部分的双栏模式中的命令当中的figure环境全部替换为figure*环境，其他命令不变；（二）在原来的命令的末尾加上多行"\blindtext"命令，将双栏中的文本内容填充到三页的篇幅。经过这两步操作后，就会发现很有意思的打印效果：首先，所有的三张图片都被打印到了第二页或者第三页上，尽管这些图片其实很早就被插入进来，尤其是第一张图片，甚至在双栏模式开始之后的第一行就被插入进来，但是其仍然在第二页才被打印出来；其次，这三张图片都被插入进了第二页或者第三页的顶部，因此，可以发现figure*环境在双栏模式下其实并不遵循htbp中的hb规则，而只遵循tp规则（在第一页上似乎只遵循p规则，否则无法解释图片为什么没有被插入第一页的顶部）；最后，这些图片被插入的位置并不处于双栏模式下，而是处于单栏模式下，类似于"\twocolumn[text]"命令可选参数的设置效果。结合以上三点，可以发现figure*环境的作用就是在双栏模式下插入【不在双栏内部浮动，而是在每一页的双栏顶部预留出的单栏区域浮动】的图片。这一环境有其存在的合理性：和在双栏模式下的figure环境相比，其可以保证大图在双栏模式下依然拥有足够的插入空间，与双栏模式之外的figure环境相比，其可以保证图片仍然可以插入到双栏文本的内部，而不是只能出现在双栏文本的前面或者后面，table*环境和table环境的区别也可以参照理解。
    %! figure*环境和table*环境中的"*"号和章节命令后面的"*"号代表了完全不同的功能，前者的功能上文已经详述，后者则会导致计数器停止计数，因此，需要意识到，LaTeX中的"*"号并不是一个带有专一功能的记号，而是一个代表在原有命令/环境的基础上改变某些功能的记号
    
%rfr 关于上述figure*环境的一些不足之处，多个宏包提出了解决的办法
%rfr （三）首先是关于figure*环境中的图片在双栏模式下只能在下一页被插入的问题，可以使用cuted宏包提供的strip环境来解决，可以参考：https://latex.org/forum/viewtopic.php?f=45&t=10661
\twocolumn
    % 以下strip环境由cuted宏包提供
    \begin{strip}
        \centering
        \includegraphics[width = .1\textwidth]{doraemon1.jpg}
        \captionof{figure}{这是用strip环境在双栏模式下插入的图片} %! 在strip环境中需要使用caption宏包提供的"\captionof{float type}{heading}"命令来设置标题
    \end{strip}

    \textbf{第三部分：strip环境}

    \blindtext[2]

%rfr （四）关于figure*环境中的图片在双栏模式下只能在双栏顶部被插入的问题，可以使用dblfloatfix宏包来解决（"dbl"来自"double (column)"的缩写），但是仍然无法解决图片无法在首页被插入的问题，可以参考：https://tex.stackexchange.com/questions/167186/figure-environment-skips-page-while-using-two-column-document
%! 由于dblfloatfix宏包和stfloats宏包不兼容，因此在此文档中没有加载dblfloatsfix宏包，打印出的实际上是stfloats宏包的效果，如果想要查看dblflatsfix宏包的效果，可以选择加载该宏包而不加载另外两个宏包
\twocolumn
    \begin{figure*}[b] % 强制设置插入位置为底部
        \includegraphics[width = .1\textwidth]{doraemon1.jpg}
        \caption{这是加载了dblfloatfix宏包后，用figure*环境在双栏模式下插入的图片}
    \end{figure*}

    \textbf{第四部分：dblfloatfix宏包}

    \blindtext[4]

% rfr （五）有两个宏包可以一次性解决上述两个问题：stfloats宏包（"st"似乎是来自"standard"的缩写）和nidanfloat宏包（这是一个日本人开发的宏包，"nidan"来自日语“二段”的罗马字，“二段”即指“双栏”），可以分别参考：https://blog.csdn.net/zhuang19951231/article/details/79176298、https://tex.stackexchange.com/questions/167186/figure-environment-skips-page-while-using-two-column-document
%! 由于nidanfloat宏包会导致通过cuted宏包提供的strip环境中的图片消失，但编译不会报错，因此此文档中只加载stfloats宏包，而不加载nidanfloat宏包，如果想要查看nidanfloat宏包的效果，可以加载该宏包，取消加载stfloats宏包
\twocolumn
    \begin{figure*}[bp]
        % 当设置为h时，图片会插入到第二页的顶部，说明h是无效设置，事实上，从figure*环境的功能出发也很容易理解这一点，h设置完全可以由figure环境来承担
        % 当设置为t(b)(p)时，图片会插入首页的顶部，这相当于cuted宏包提供的strip环境的效果
        % 当设置为b(p)时，图片会插入首页的底部，这是dblfloatfix宏包无法达到的效果
        \includegraphics[width = .1\textwidth]{doraemon1.jpg}
        \caption{这是加载了stfloats/nidanfloat宏包后，用figure*环境在双栏模式下插入的图片（一）}
    \end{figure*}

    \textbf{第五部分：stfloats宏包}

    \blindtext[3]

    \begin{figure*}[tbp]
        \includegraphics[width = .1\textwidth]{doraemon1.jpg}
        \caption{这是加载了stfloats/nidanfloat宏包后，用figure*环境在双栏模式下插入的图片（二）}
    \end{figure*}
    
    \blindtext[3]

\onecolumn

%* （六）下面讨论在multicol宏包提供的multicols环境中插入图片的问题
%! 在multicols环境中，上述的很多命令都会失效：
%! （1）使用figure环境插入图片，虽然编译不会报错，但是图片实际不会被插入进多栏中，使用cuted宏包提供的strip环境插入图片，也会产生类似的效果，在multocols环境中，只能使用figure*环境插入图片
%! （2）加载nidanfloat宏包会导致multicols环境中使用figure*环境插入的图片消失，但编译不会报错，类似于其对strip环境中的图片产生的影响
% （3）加载stfloats宏包，能够顺利使用figure*环境在多栏中插入图片，但是图片无法插入到首页当中，而只能从第二页开始插入
% （4）加载dblfloatfix宏包，产生的也是和（3）中一样的效果，这倒是和dblfloatfix宏包本身提供的功能相符合
%* （5）事实上，就算不加载stfloats宏包和dbfloatfix宏包，（3）和（4）中的效果仍然存在，说明这是multicols环境自带的效果，具体情况请看下文注释

        \begin{figure}[t]
            \includegraphics[width = .1\textwidth]{doraemon1.jpg}
            \caption{这是在multicols环境外插入的图片（首页顶部）}
        \end{figure}

        \begin{figure}[b]
            \includegraphics[width = .1\textwidth]{doraemon1.jpg}
            \caption{这是在multicols环境外插入的图片（首页底部）}
        \end{figure}
        %rfr 如果要达到在多栏的首页插入图片的效果，只能在multicols环境外之前的位置插入图片（此时选择figure环境还是figure*环境都可以），如果设置t(b)(p)，则插入首页的顶部，如果设置b(p)，则插入首页的底部，可以参考：https://latex.org/forum/viewtopic.php?f=45&t=10661
\begin{multicols}{3}[\centering 这是分成三栏的前言] 
        % 和"\twocolumn[text]"命令的可选参数一样，相应内容会打印在多栏顶部预留出的单栏区域

        \begin{figure*}[t]
            %* 当设置为t(b)(p)时，图片会从第二页开始插入合适页面的顶部，这也是默认情况下的效果
            \includegraphics[width = .1\textwidth]{doraemon1.jpg}
            \caption{这是在multicols环境中用figure*环境插入的图片（首页以后的顶部）}
        \end{figure*}

        \textbf{第六部分：multicols环境}

        \blindtext[2]

        \begin{figure*}[b]
            %* 当设置为b(p)时，图片会从第二页开始插入合适页面的底部，这相当于dblfloatfix宏包提供的效果
            \includegraphics[width = .1\textwidth]{doraemon1.jpg}
            \caption{这是在multicols环境中用figure*环境插入的图片（首页以后的底部）}
        \end{figure*}

        \columnbreak

        \mbox{}

        \columnbreak

        \blindtext[2]

\end{multicols}


\blindtext

\end{document}